% Hlavička - nastavení dokumentu

\documentclass[a5paper,
		       11pt]{article} % nastaveni tak aby bylo pismo dostatecne velike
\usepackage[a5paper,
			top = 0.7cm,
			left = 1cm,
			right = 1cm,
			bottom = 1.5cm]{geometry} % Knihovna na nastavení velikosti okrajů stránky
\usepackage [utf8] {inputenc} % Kódování textu. Píšeme v utf-8, aby to pěkně jelo i linuxákům a všem a nedělalo to čínský znaky
\usepackage[T1]{fontenc}
\usepackage{babel} % nějaké další potřebné balíčky
\usepackage{guitar} %knihovna pro sazbu akordů
\usepackage{multicol} %knihovna pro možnost sazby ve více sloupcích
\usepackage{graphicx} % knihovna pro vkládání obrázků
\usepackage{setspace}

% následuje definice maker pro přehlednější práci

\newcommand{\ch}[1]{\guitarChord{\textbf{#1}}}
\newcommand{\song}[1]{\newpage \subsection{#1} \begin{enumerate}} 
\newcommand{\esong}{\end{enumerate}}
\newcommand{\ve}{\item }				% sloka
\newcommand{\cho}[1]{\item[R:]#1}		% refrén
\newcommand{\re}[1]{\item[#1]}			% něco jiného, odstavec
\newcommand{\add}[1]{\input{songs/#1.tex}}

%widowpenalties 1 10000 %zakáže zlom stránky uprostřed sloky

\author{Mara a Martin} % je zvykem vyplnit
\title{Zpěvník -- nástroj} % taky

\setcounter{secnumdepth}{1} % Nastavení, aby se nečíslovaly názvy písniček jako kapitoly

\widowpenalties 1 10000 				% zakáže zlom stránky uprostřed sloky
\raggedbottom

% ZAČÁTEK SAMOTNÉHO DOKUMENTU
\begin{document}

\song{Rád chodím na poštu (Pokáč)}

\cho{}
\ch{D}Rád chodím \ch{A}na poštu,\\
\ch{G}nejsem tam jen \ch{D}do počtu\\
\ch{G}jsem tam důležitý \ch{D}člen\\
\ch{Emi}fronty která tr\ch{A}{č}í ven

vždy když se cítím sám,\\
rád na poštu zavítám\\
tam je tolik lidí, že\\
\ch{Emi}pohnout se má\ch{A}m potíž\ch{D}e

\ve{}
\ch{G}za přepážkou \ch{D}paní v letech \ch{G}má razítko \ch{D}gumové\\
\ch{G}a občas tiskne \ch{Hmi}tiskárnou z dob \ch{Emi}první války \ch{A}světové\\
každému kdo spílá jí že zas na poště pros*al den\\
nabídne los stírací ať uklidní se hazardem

\cho{}

\ve{}
proč je vždy jen jedna z pěti přepážek otevřená\\
je záhada co navždy zůstat má tajemstvím zastřená\\
má nejoblíbenější služba je balíček do ruky\\
znamená to totiž že zas budu moct jít na poštu\\
\ch{A}jestli ho vůbec doručej\\
\ch{A}to už je bez záruky

\cho{}
...trčí ven\\
vždy když tam stepuju\\
tenhle song si notuju\\
mám pak hnedka trochu míň\\
chuť \ch{Emi}do mozku si \ch{A}vrazit \ch{D}klín

\esong{}



\end{document}


